%! Parametric Functions Modeling Planar Motion — Unit 4, Lesson 4.2 — Guided Notes (Answer Key)
\documentclass[10pt]{article}
\usepackage{precalculus-article}
\usepackage{precalculus-key}   % pulls in -boxes; do NOT also load -boxes
\newcommand{\TallMath}[1]{$\displaystyle #1\rule[-1.4em]{0pt}{3.2em}$}

\begin{document}

\pageheader{Unit 4, Lesson 4.2}{Guided Notes --- Answer Key}
\namedateperiod

\vspace{0.08in}

%----------------------------------------------------------------------
% Objectives
%----------------------------------------------------------------------
\begin{objectivebox}
\textbf{I will be able to\ldots}
\begin{enumerate}[leftmargin=1.5em, itemsep=4pt]
  \item \textbf{LO 4.2.A} --- Identify key characteristics of a
    parametric planar motion function related to position: find the
    horizontal and vertical extrema from the max/min values of $x(t)$
    and $y(t)$, and locate $x$- and $y$-intercepts of the curve from
    the zeros of $y(t)$ and $x(t)$. \writeline
\end{enumerate}
\end{objectivebox}

\vspace{0.07in}

%----------------------------------------------------------------------
% Vocabulary
%----------------------------------------------------------------------
\begin{vocabbox}
\begin{description}[leftmargin=0pt, itemsep=5pt]
  \item[\termblanklong{horizontal extremum}]
    The \ans{maximum} or \ans{minimum}
    value of $x(t)$; gives the rightmost or leftmost position of the
    particle.

  \item[\termblanklong{vertical extremum}]
    The maximum or minimum value of \ans{$y(t)$}; gives the
    highest or lowest position of the particle.

  \item[\termblanklong{$y$-intercept (of parametric curve)}]
    A point on the curve where $x(t) = $ \ans{$0$};
    found by solving $x(t) = 0$ for $t$, then evaluating
    $y(t)$ at those values.

  \item[\termblanklong{$x$-intercept (of parametric curve)}]
    A point on the curve where $y(t) = $ \ans{$0$};
    found by solving $y(t) = 0$ for $t$, then evaluating $x(t)$.
\end{description}
\end{vocabbox}

\vspace{0.07in}

%----------------------------------------------------------------------
% Hook
%----------------------------------------------------------------------
\begin{hookbox}
\textbf{Launch:} A ball's position (feet) at time $t$ seconds is
$x(t) = -16t^2 + 48t$ and $y(t) = 12t - t^2$.

\vspace{4pt}
\textbf{Q1.} Which component do you examine to find how high the ball
goes? What equation do you solve to find when the ball lands?

\vspace{4pt}
\ans{Examine $y(t)$ for height. To find when the ball lands, solve
$y(t) = 12t - t^2 = 0$.}
\end{hookbox}

\vspace{0.07in}

%----------------------------------------------------------------------
% LO 4.2.A — Extrema
%----------------------------------------------------------------------
\begin{notesbox}{LO 4.2.A --- Planar Motion: Horizontal and Vertical Extrema}

\textbf{Example:} Let $f(t) = (x(t),\, y(t)) = (t^2 - 4t + 3,\;
{-t^2 + 2t})$ for $0 \le t \le 4$.

\vspace{6pt}
\textbf{Step 1 --- Build the table.}

\vspace{4pt}
\begin{center}
\renewcommand{\arraystretch}{1.6}
\begin{tabular}{|c|c|c|c|}
\hline
$t$ & $x(t) = t^2 - 4t + 3$ & $y(t) = -t^2 + 2t$ & $(x,\; y)$ \\
\hline
$0$ & \ans{$3$}  & \ans{$0$}  & \ans{$(3,\;0)$}  \\
\hline
$1$ & \ans{$0$}  & \ans{$1$}  & \ans{$(0,\;1)$}  \\
\hline
$2$ & \ans{$-1$} & \ans{$0$}  & \ans{$(-1,\;0)$} \\
\hline
$3$ & \ans{$0$}  & \ans{$-3$} & \ans{$(0,\;-3)$} \\
\hline
$4$ & \ans{$3$}  & \ans{$-8$} & \ans{$(3,\;-8)$} \\
\hline
\end{tabular}
\end{center}

\vspace{8pt}
\textbf{Step 2 --- Find horizontal extrema of $x(t) = t^2 - 4t + 3$.}

$x(t)$ is a \ans{upward-opening} parabola. Vertex at
$t = -\dfrac{b}{2a} = -\dfrac{-4}{2} = $ \ans{$2$}

Minimum value of $x$: $x(2) = $ \ans{$-1$}
\quad (particle's \ans{left}-most position)

Maximum value of $x$ on $[0,4]$: check endpoints.
$x(0) = $ \ans{$3$}, \quad $x(4) = $ \ans{$3$}.
Maximum is \ans{$3$} (particle's \ans{right}-most position).

\vspace{8pt}
\textbf{Step 3 --- Find vertical extrema of $y(t) = -t^2 + 2t$.}

Vertex at $t = -\dfrac{b}{2a} = $ \ans{$1$}.

Maximum value of $y$: $y(1) = $ \ans{$1$}
\quad (particle's \ans{high}-est position)

Minimum of $y$ on $[0,4]$: check endpoints.
$y(0) = $ \ans{$0$}, \quad $y(4) = $ \ans{$-8$}.
Minimum is \ans{$-8$}.

\end{notesbox}

\vspace{0.07in}

%----------------------------------------------------------------------
% LO 4.2.A — Intercepts
%----------------------------------------------------------------------
\begin{notesbox}{LO 4.2.A --- Axis Intercepts of a Parametric Curve}

\textbf{Key rule:} Zeros of $x(t)$ give \ans{$y$}-intercepts
of the curve. Zeros of $y(t)$ give \ans{$x$}-intercepts.

\vspace{8pt}
\textbf{Example (continued):} $f(t) = (t^2 - 4t + 3,\; {-t^2 + 2t})$,
$0 \le t \le 4$.

\vspace{4pt}
\textbf{Find the $y$-intercepts} (where $x(t) = 0$):

Set $t^2 - 4t + 3 = 0$. Factor: $(t - $ \ans{$1$}$)(t - $
\ans{$3$}$) = 0$.

$t = $ \ans{$1$} or $t = $ \ans{$3$}.

Evaluate $y$ at each: $y(1) = $ \ans{$1$},
\quad $y(3) = $ \ans{$-3$}.

$y$-intercepts of the curve: \ans{$(0,\,1)$} and \ans{$(0,\,-3)$}.

\vspace{8pt}
\textbf{Find the $x$-intercepts} (where $y(t) = 0$):

Set $-t^2 + 2t = 0$. Factor: $-t(t - $ \ans{$2$}$) = 0$.

$t = $ \ans{$0$} or $t = $ \ans{$2$}.

Evaluate $x$ at each: $x(0) = $ \ans{$3$},
\quad $x(2) = $ \ans{$-1$}.

$x$-intercepts of the curve: \ans{$(3,\,0)$} and \ans{$(-1,\,0)$}.

\end{notesbox}

\end{document}
